The support fraction divides distinct observed coordinates by the number of candidate coordinates in the current table. The distance features use all support-range-normalized coordinates. Local span is the Euclidean norm of the coordinate ranges of the enclosing vertices, zero outside. Envelope width is divided component-wise by the source response scale. Distance, span and envelope width use $\log(1+z)$. Typical spacing is the median nearest-neighbor distance among distinct observations. A positive typical spacing and a local span at most twice that spacing identify dense support. Outside membership takes precedence, followed by an unobserved width strictly between the minimum and maximum observed widths, then the dense rank-one and higher-rank cases.

For variation and slopes, order the three nearest observations by distance with stable input-order ties. Variation is the component range of their scaled responses. The first slope divides the second-minus-first response by the distance between those observations, floored at $10^{-12}$. Curvature subtracts that slope from the corresponding second-to-third slope. Unavailable derivatives are zero. The gate receives log typical spacing and variation, signed-log slopes and curvature, and the same context as the geometry tokens. Material tokens occupy the same feature width as geometry tokens, $17+2K$.
